\documentclass{article}
\usepackage{amsmath}
\usepackage{hyperref}
\usepackage{graphicx}

\title{My First Document}
\author{Jane Doe}
\date{\today}

\begin{document}

\maketitle

% ── Introduction ──────────────────────────────────────────────────────────────

\section{Introduction}

Welcome to the \textbf{LaTeX Preview Editor}. This document exercises
all currently supported decorations. Text can be \textit{italic},
\textbf{bold}, \underline{underlined}, or in \texttt{monospace}.
You can also \emph{emphasize} words for semantic meaning.

Cross-references look like chips inline: see \ref{sec:math} for equations
and \ref{sec:lists} for list examples. Citations render similarly:
\cite{einstein1905} and \cite{dirac1930}.

\label{sec:intro}

% ── Mathematics ───────────────────────────────────────────────────────────────

\section{Mathematics}
\label{sec:math}

Inline math: the energy $E = mc^2$ is famous, and so is Euler's
identity $e^{i\pi} + 1 = 0$.

A display equation using \verb|\[...\]|:

\[
  \int_0^\infty e^{-x^2} \, dx = \frac{\sqrt{\pi}}{2}
\]

A numbered equation:

\begin{equation}
  \nabla \cdot \mathbf{E} = \frac{\rho}{\varepsilon_0}
  \label{eq:gauss}
\end{equation}

See \eqref{eq:gauss} for Gauss's law.  An aligned system:

\begin{align}
  f(x) &= x^2 + 2x + 1 \\
       &= (x + 1)^2
\end{align}

% ── Lists ─────────────────────────────────────────────────────────────────────

\section{Lists}
\label{sec:lists}

An unordered list:

\begin{itemize}
  \item First bullet point with some text
  \item Second bullet, containing inline math $\alpha + \beta$
  \item Third bullet point
\end{itemize}

A numbered list:

\begin{enumerate}
  \item Step one: set up the environment
  \item Step two: write the content
  \item Step three: compile and review
\end{enumerate}

% ── Figures and Tables ────────────────────────────────────────────────────────

\section{Figures and Tables}
\label{sec:figures}

Here is an example figure environment:

\begin{figure}
  \centering
  \includegraphics[width=0.5\textwidth]{example-image}
  \caption{An example figure showing LaTeX rendering}
  \label{fig:example}
\end{figure}

See Figure~\ref{fig:example} for details. Here is a table example:

\begin{table}
  \centering
  \begin{tabular}{|c|c|c|}
    \hline
    Column 1 & Column 2 & Column 3 \\
    \hline
    Data 1 & Data 2 & Data 3 \\
    \hline
  \end{tabular}
  \caption{A sample data table}
  \label{tab:sample}
\end{table}

See Table~\ref{tab:sample} for the data.

% ── URLs and footnotes ────────────────────────────────────────────────────────

\section{Links and Footnotes}

Visit \url{https://www.overleaf.com} for online LaTeX editing.
This sentence has a footnote.\footnote{Footnotes appear inline as small
dimmed text until you click on them.}

\end{document}
